\begin{definition}[Actual open graph coefficients]%
    \label{def:peps_graph_open_region_contraction}
    \lean{TNLean.PEPS.graphRegionIncidentConfigEquiv,
        TNLean.PEPS.graphOpenRegionNetwork}
    \leanok
    Let $R$ be a region in a finite simple graph. Incident virtual labels
    split into internal-bond labels and crossing-bond labels. For a finite
    virtual alphabet $X$, site coefficients $A_v$, prescribed crossing
    labels $\theta$, and physical configuration $\sigma$, define the actual
    open coefficient by summing the product of the site coefficients over
    incident assignments whose crossing restriction equals $\theta$.
    The physical alphabet may depend on $v$.
    These are the finite contraction coordinates used in
    \cite[Sections~6.2--6.3 and~7]{Schuch2010PEPS}.
\end{definition}

\begin{theorem}[Exact internal summation and numbered open coefficients]%
    \label{thm:peps_graph_open_region_contraction}
    \lean{TNLean.PEPS.graphRegionIncidentConfigEquiv_apply_internal,
        TNLean.PEPS.graphRegionIncidentConfigEquiv_apply_boundary,
        TNLean.PEPS.graphOpenRegionNetwork_eq_sum_internal,
        TNLean.PEPS.openRegionWeight_groupBondTensor_eq_graphOpenRegionNetwork}
    \leanok
    \uses{def:peps_graph_open_region_contraction,def:peps_open_region_contraction}
    The two factors of the incident-label bijection read exactly the original
    internal and crossing labels. Prescribing $\theta$ therefore gives
    \begin{align}
        A_R(\theta,\sigma)
        =\sum_{\xi\text{ internal}}\prod_{v\in R}
            A_v\bigl((\xi,\theta)|_v,\sigma_v\bigr),
    \end{align}
    with each internal assignment counted once. When the physical alphabet
    is uniformly numbered, enumerating $X$ into finite indices gives exactly
    the original numbered open-region coefficient of the same site tensors.
    No boundary-state equality or Gram identity is a hypothesis.
\end{theorem}

\begin{proof}\leanok
    Partition incident edges into internal and crossing edges and transport
    their labels along the resulting bijection. The prescribed crossing
    condition removes the crossing-label sum. For numbered coefficients,
    reindex each incident virtual label by the same finite enumeration;
    the site products and boundary condition are unchanged.
\end{proof}
